\let\negmedspace\undefined
\let\negthickspace\undefined
\documentclass[journal]{IEEEtran}
\usepackage[a5paper, margin=10mm, onecolumn]{geometry}
%\usepackage{lmodern} % Ensure lmodern is loaded for pdflatex
\usepackage{tfrupee} % Include tfrupee package

\setlength{\headheight}{1cm} % Set the height of the header box
\setlength{\headsep}{0mm}     % Set the distance between the header box and the top of the text

\usepackage{gvv-book}
\usepackage{gvv}
\usepackage{cite}
\usepackage{amsmath,amssymb,amsfonts,amsthm}
\usepackage{algorithmic}
\usepackage{graphicx}
\usepackage{textcomp}
\usepackage{xcolor}
\usepackage{txfonts}
\usepackage{listings}
\usepackage{enumitem}
\usepackage{mathtools}
\usepackage{gensymb}
\usepackage{comment}
\usepackage[breaklinks=true]{hyperref}
\usepackage{tkz-euclide} 
\usepackage{listings}
% \usepackage{gvv}                               

\def\inputGnumericTable{}                      
\usepackage[latin1]{inputenc}                    
\usepackage{color}                              
\usepackage{array}                             
\usepackage{longtable}                          
\usepackage{calc}                               
\usepackage{multirow}                           
\usepackage{hhline}                            
\usepackage{ifthen}                          
\usepackage{lscape}
\begin{document}

\bibliographystyle{IEEEtran}
\vspace{3cm}

\title{7.2.23}
\author{AI25BTECH11024 - Pratyush Panda
}
\maketitle
% \newpage
% \bigskip
{\let\newpage\relax\maketitle}

\renewcommand{\thefigure}{\theenumi}
\renewcommand{\thetable}{\theenumi}
\setlength{\intextsep}{10pt} % Space between text and floats


\numberwithin{equation}{enumi}
\numberwithin{figure}{enumi}
\renewcommand{\thetable}{\theenumi}

\textbf{Question: } \\
The area of the circle centred at $\brak{1,2}$ and passing through $\brak{4,6}$ is
\begin{enumerate}
\item $5\pi$
\item $10\pi$
\item $25\pi$
\item None of the above
\end{enumerate}
\vspace{0.7cm}

\renewcommand{\theequation}{0.\arabic{equation}}

\textbf{Solution: } \\
The centre of the circle is given to us as $\Vec{c}=\myvec{1 \\ 2}$.
Thus we can write the $\Vec{u}$ as:
\begin{align}
\Vec{u}=-\Vec{c}=\myvec{-1 \\ -2}
\end{align}

We also know that $f=||\Vec{c}||^2-r^2$, where $r$ is the radius of the circle.\\

To find the area of the circle we need to find the radius of the circle. For that we can use the equation of circle.\\
The equation of the circle can be written as;
\begin{align}
||\Vec{x}||^2+2\Vec{u}^T\Vec{x}+f=0 \\
or, \, ||\Vec{x}||^2+2\Vec{u}^T\Vec{x}+||\Vec{c}||^2-r^2=0 \\
\implies r^2=||\Vec{x}||^2+2\Vec{u}^T\Vec{x}+||\Vec{c}||^2 
\end{align}
where $\Vec{x}$ is any point on the circle \\
Since we are given that $\brak{4,6}$ is a point on the circle, we can put $\Vec{x}=\myvec{4 \\ 6}$

On putting the values in the equation $\brak{0.3}$, we get;
\begin{align}
r^2=52-32+5=25
\end{align}
Thus, the area of the circle$\brak{A}$ will be:
\begin{align}
A=\pi r^2=25\pi
\end{align}

\begin{figure}[H]
\centering
\includegraphics[width=0.8\columnwidth]{figs/img.png}
\caption*{}
\end{figure}

\end{document}
